\subsection{ECHELONS}

An echelon construction scheme is a sequence \(c_1, c_2, \ldots, c_m\) of ordered pairs of natural integers (*) \(c_i = (a_i, b_i)\) , satisfying the following conditions:

\begin{enumerate}
\def\labelenumi{(\alph{enumi})}
\item
  If \(b_{i} = 0\) , then \(1 \leqslant a_{i} \leqslant i - 1\) .
\item
  If \(a_i \neq 0\) and \(b_i \neq 0\) , then \(1 \leqslant a_i \leqslant i - 1\) and \(1 \leqslant b_i \leqslant i - 1\) .
\end{enumerate}

These conditions imply that \(c_{1}=(0,b_{1})\) , with \(b_{1}>0\) . If n is the largest of the integers \(b_{i}\) which appear in the pairs \((0,b_{i})\) , then \(c_{1},c_{2},\ldots,c_{m}\) is said to be an echelon construction scheme on n terms.

Given an echelon construction scheme \(\mathbf{S} = (c_{1}, c_{2}, \ldots, c_{m})\) on n terms, and given n terms \(E_{1}, E_{2}, \ldots, E_{n}\) in a theory C which is stronger than the theory of sets, an echelon construction of scheme S on \(E_{1}, \ldots, E_{n}\) is defined to be a sequence \(A_{1}, A_{2}, \ldots, A_{m}\) of m terms in the theory C, defined step by step by the following conditions:

\begin{enumerate}
\def\labelenumi{(\alph{enumi})}
\item
  If \(c_{i} = (0, b_{i})\) , then \(A_{i}\) is the term \(E_{b_{i}}\) .
\item
  If \(c_{i} = (a_{i}, 0)\) , then \(A_{i}\) is the term \(\mathfrak{P}(A_{a_i})\) .
\item
  If \(c_{i} = (a_{i}, b_{i})\) , where \(a_{i} \neq 0\) and \(b_{i} \neq 0\) , then \(A_{i}\) is the term \(A_{a_{i}} \times A_{b_{i}}\) .
\end{enumerate}

The last term \(A_{m}\) of the echelon construction of scheme S on \(E_{1}, \ldots, E_{n}\) is called the echelon of scheme S on the base sets \(E_{1}, \ldots, E_{n}\) ; in the general arguments which follow, it will be denoted by the notation \(S(E_{1}, \ldots, E_{n})\) .

Example. Given two sets E, F, the set \(\mathfrak{P}(\mathfrak{P}(E)) \times \mathfrak{P}(F)\) is an echelon on E, F, with scheme

\[
(0, 1), (0, 2), (1, 0), (3, 0), (2, 0), (4, 5).
\]

It is also the echelon on E, F with scheme

\[
(0, 2), (0, 1), (1, 0), (2, 0), (4, 0), (5, 3).
\]

Distinct schemes may therefore give rise to the same echelon on the same terms.

